
\documentclass{article}
\usepackage{colortbl}
\usepackage{makecell}
\usepackage{multirow}
\usepackage{supertabular}

\begin{document}

\newcounter{utterance}

\centering \large Interaction Transcript for game `clean\_up', experiment `0\_easy\_3obj\_lv', episode 0 with Apertus{-}v1.5{-}8B.
\vspace{24pt}

{ \footnotesize  \setcounter{utterance}{1}
\setlength{\tabcolsep}{0pt}
\begin{supertabular}{c@{$\;$}|p{.15\linewidth}@{}p{.15\linewidth}p{.15\linewidth}p{.15\linewidth}p{.15\linewidth}p{.15\linewidth}}
    \# & $\;$A & \multicolumn{4}{c}{Game Master} & $\;\:$B\\
    \hline

    \theutterance \stepcounter{utterance}  
    & & \multicolumn{4}{p{0.6\linewidth}}{
        \cellcolor[rgb]{0.9,0.9,0.9}{
            \makecell[{{p{\linewidth}}}]{
                \texttt{\tiny{[P1$\langle$GM]}}
                \texttt{Es esmu jūsu spēles vadītājs, un jūs spēlējat sadarbības spēli ar šādu režģi kā spēles laukumu:} \\
\\ 
\texttt{```} \\
\texttt{ 1234567} \\
\texttt{╔═╤═══╤═╗} \\
\texttt{║◌│◌◌◌│◌║ 1} \\
\texttt{╟─┼──┬┘◌║ 2} \\
\texttt{║◌│◌◌│◌P║ 3} \\
\texttt{╟─┴──┘◌◌║ 4} \\
\texttt{║◌L◌◌◌◌◌║ 5} \\
\texttt{║◌C◌◌◌◌◌║ 6} \\
\texttt{║◌◌◌◌◌◌◌║ 7} \\
\texttt{╚═══════╝} \\
\\ 
\texttt{```} \\
\\ 
\texttt{* Augšējā malā ir attēlotas x koordinātas, kas pieaug pa labi, un labajā malā y koordinātas, kas pieaug uz leju.} \\
\texttt{* Uz jūsu režģa ir nejauši izvietoti šādi objekti: 'C', 'L', 'P'.} \\
\\ 
\\ 
\texttt{Otrs spēlētājs redz spēles laukuma variāciju, kurā objekti ir izvietoti citās nejaušās vietās. Jūs nevarat redzēt otra spēlētāja laukumu, un viņš nevar redzēt jūsējo.} \\
\\ 
\\ 
\texttt{**Mērķis:**} \\
\\ 
\texttt{Abiem spēlētājiem ir jāpārvieto objekti uz sava fona tā, lai vienādi objekti nonāktu vienādās koordinātēs. Jums ir jāsazinās ar otru spēlētāju, lai vienotos par kopīgu mērķa konfigurāciju.} \\
\\ 
\\ 
\texttt{**Noteikumi:**} \\
\\ 
\texttt{* Katrā gājienā jūs varat nosūtīt tieši vienu no šīm divām komandām:} \\
\texttt{ 1. `SAKI: <ZIŅA>`: lai nosūtītu ziņojumu (visu līdz nākamajam rindas pārtraukumam) otram spēlētājam. Es to pārsūtīšu jūsu partnerim.} \\
\texttt{2. `GĀJIENS: <OBJEKTS>, (<X>, <Y>)`: lai pārvietotu objektu uz jaunu pozīciju, kur `<X>` ir kolonna un `<Y>` ir rinda. Es jūs informēšu, vai jūsu gājiens bija veiksmīgs vai nē.} \\
\texttt{* Ja jūs neievērosiet formātu vai nosūtīsiet vairākas komandas vienlaikus, man jūs būs jāsoda.} \\
\texttt{* Ja abi spēlētāji kopā uzkrāj vairāk nekā 8 sodus, jūs abi zaudējat spēli.} \\
\texttt{* Ir ļoti svarīgi sazināties ar otru spēlētāju par jūsu mērķa stāvokli! *Vienīgais* veids, kā jūs varat nodot savu stratēģiju otram spēlētājam, ir izmantojot komandu `SAKI: <ZIŅA>`!} \\
\\ 
\\ 
\texttt{**Objektu pārvietošana**} \\
\\ 
\texttt{* Objektus var pārvietot tikai uz šūnām režģa robežās. Mērķa šūnai jābūt tukšai, t.i., tajā jābūt tikai simbolam '◌'.} \\
\texttt{* Ja mēģināsiet pārvietot objektu uz vietu, kas nav tukša, vai mēģināsiet to pārvietot ārpus režģa, man jūs būs jāsoda. Jūs saņemat vēl vienu mēģinājumu.} \\
\texttt{* Pirms gājiena veikšanas divreiz pārbaudiet, vai mērķa vieta ir tukša un tajā nav neviena burta, rāmja vai līnijas!} \\
\\ 
\\ 
\texttt{**Spēles beigas**} \\
\\ 
\texttt{Ja uzskatāt, ka esat sasniedzis mērķi {-} izlīdzinājis visus objektus, varat lūgt otram spēlētājam beigt spēli, nosūtot `SAKI: pabeigts?`. Ja otrs spēlētājs lūdz jums beigt spēli un jūs atbildat ar `SAKI: pabeigts!`, spēle beigsies.} \\
\\ 
\texttt{Abi spēlētāji uzvar, ja spēle beidzas 12 raundos, kur viens raunds ir definēts kā divi spēlētāji, kas katrs nosūta derīgu komandu.} \\
\\ 
\texttt{**Punktu skaitīšana:**} \\
\\ 
\texttt{Jo tuvāk vienādi objekti atrodas abos spēles laukumos, jo vairāk punktu jūs saņemat. Sodi samazina jūsu punktus. Vai varat pārspēt rekordu?} \\
\\ 
\\ 
\texttt{Lūdzu, nosūtiet ziņu otram spēlētājam, lai sāktu spēli!} \\
            }
        }
    }
    & \\ \\

    \theutterance \stepcounter{utterance}  
    & \multicolumn{4}{p{0.6\linewidth}}{
        \cellcolor[rgb]{0.9,0.9,0.9}{
            \makecell[{{p{\linewidth}}}]{
                \texttt{\tiny{[P1$\rangle$GM]}}
                \texttt{SAKI: Esmu gatavs sākt sadarbības spēli. Lūdzu, nosūtiet man jūsu pašreizējo objektu izvietojumu, lai varētu plānot mērķa pozīcijas.} \\
            }
        }
    }
    & & \\ \\

    \theutterance \stepcounter{utterance}  
    & & & \multicolumn{2}{p{0.3\linewidth}}{
        \cellcolor[rgb]{0.9,0.9,0.9}{
            \makecell[{{p{\linewidth}}}]{
                \texttt{\tiny{[GM$|$GM]}}
                \texttt{SAKI: Esmu gatavs sākt sadarbības spēli. Lūdzu, nosūtiet man jūsu pašreizējo objektu izvietojumu, lai varētu plānot mērķa pozīcijas.} \\
            }
        }
    }
    & & \\ \\

    \theutterance \stepcounter{utterance}  
    & & \multicolumn{4}{p{0.6\linewidth}}{
        \cellcolor[rgb]{0.9,0.9,0.9}{
            \makecell[{{p{\linewidth}}}]{
                \texttt{\tiny{[GM$\rangle$P2]}}
                \texttt{Es esmu jūsu spēles vadītājs, un jūs spēlējat sadarbības spēli ar šādu režģi kā spēles laukumu:} \\
\\ 
\texttt{```} \\
\texttt{ 1234567} \\
\texttt{╔═╤═══╤═╗} \\
\texttt{║◌│◌C◌│◌║ 1} \\
\texttt{╟─┼──┬┘◌║ 2} \\
\texttt{║◌│◌◌│◌◌║ 3} \\
\texttt{╟─┴──┘◌◌║ 4} \\
\texttt{║◌◌◌◌◌◌◌║ 5} \\
\texttt{║◌◌◌◌◌◌◌║ 6} \\
\texttt{║◌◌◌P◌◌L║ 7} \\
\texttt{╚═══════╝} \\
\\ 
\texttt{```} \\
\\ 
\texttt{* Augšējā malā ir attēlotas x koordinātas, kas pieaug pa labi, un labajā malā y koordinātas, kas pieaug uz leju.} \\
\texttt{* Uz jūsu režģa ir nejauši izvietoti šādi objekti: 'C', 'L', 'P'.} \\
\\ 
\\ 
\texttt{Otrs spēlētājs redz spēles laukuma variāciju, kurā objekti ir izvietoti citās nejaušās vietās. Jūs nevarat redzēt otra spēlētāja laukumu, un viņš nevar redzēt jūsējo.} \\
\\ 
\\ 
\texttt{**Mērķis:**} \\
\\ 
\texttt{Abiem spēlētājiem ir jāpārvieto objekti uz sava fona tā, lai vienādi objekti nonāktu vienādās koordinātēs. Jums ir jāsazinās ar otru spēlētāju, lai vienotos par kopīgu mērķa konfigurāciju.} \\
\\ 
\\ 
\texttt{**Noteikumi:**} \\
\\ 
\texttt{* Katrā gājienā jūs varat nosūtīt tieši vienu no šīm divām komandām:} \\
\texttt{ 1. `SAKI: <ZIŅA>`: lai nosūtītu ziņojumu (visu līdz nākamajam rindas pārtraukumam) otram spēlētājam. Es to pārsūtīšu jūsu partnerim.} \\
\texttt{2. `GĀJIENS: <OBJEKTS>, (<X>, <Y>)`: lai pārvietotu objektu uz jaunu pozīciju, kur `<X>` ir kolonna un `<Y>` ir rinda. Es jūs informēšu, vai jūsu gājiens bija veiksmīgs vai nē.} \\
\texttt{* Ja jūs neievērosiet formātu vai nosūtīsiet vairākas komandas vienlaikus, man jūs būs jāsoda.} \\
\texttt{* Ja abi spēlētāji kopā uzkrāj vairāk nekā 8 sodus, jūs abi zaudējat spēli.} \\
\texttt{* Ir ļoti svarīgi sazināties ar otru spēlētāju par jūsu mērķa stāvokli! *Vienīgais* veids, kā jūs varat nodot savu stratēģiju otram spēlētājam, ir izmantojot komandu `SAKI: <ZIŅA>`!} \\
\\ 
\\ 
\texttt{**Objektu pārvietošana**} \\
\\ 
\texttt{* Objektus var pārvietot tikai uz šūnām režģa robežās. Mērķa šūnai jābūt tukšai, t.i., tajā jābūt tikai simbolam '◌'.} \\
\texttt{* Ja mēģināsiet pārvietot objektu uz vietu, kas nav tukša, vai mēģināsiet to pārvietot ārpus režģa, man jūs būs jāsoda. Jūs saņemat vēl vienu mēģinājumu.} \\
\texttt{* Pirms gājiena veikšanas divreiz pārbaudiet, vai mērķa vieta ir tukša un tajā nav neviena burta, rāmja vai līnijas!} \\
\\ 
\\ 
\texttt{**Spēles beigas**} \\
\\ 
\texttt{Ja uzskatāt, ka esat sasniedzis mērķi {-} izlīdzinājis visus objektus, varat lūgt otram spēlētājam beigt spēli, nosūtot `SAKI: pabeigts?`. Ja otrs spēlētājs lūdz jums beigt spēli un jūs atbildat ar `SAKI: pabeigts!`, spēle beigsies.} \\
\\ 
\texttt{Abi spēlētāji uzvar, ja spēle beidzas 12 raundos, kur viens raunds ir definēts kā divi spēlētāji, kas katrs nosūta derīgu komandu.} \\
\\ 
\texttt{**Punktu skaitīšana:**} \\
\\ 
\texttt{Jo tuvāk vienādi objekti atrodas abos spēles laukumos, jo vairāk punktu jūs saņemat. Sodi samazina jūsu punktus. Vai varat pārspēt rekordu?} \\
\\ 
\texttt{Otrs spēlētājs sāka spēli, nosūtot šo ziņojumu:} \\
\\ 
\texttt{"Esmu gatavs sākt sadarbības spēli. Lūdzu, nosūtiet man jūsu pašreizējo objektu izvietojumu, lai varētu plānot mērķa pozīcijas."} \\
\\ 
\texttt{Kāda ir jūsu pirmā komanda?} \\
            }
        }
    }
    & \\ \\

    \theutterance \stepcounter{utterance}  
    & & & \multicolumn{4}{p{0.6\linewidth}}{
        \cellcolor[rgb]{0.9,0.9,0.9}{
            \makecell[{{p{\linewidth}}}]{
                \texttt{\tiny{[GM$\langle$P2]}}
                \texttt{SAKI: Man pašlaik ir šāds objektu izvietojums: C ir (1,2), L ir (7,5), P ir (7,6). Lūdzu, nosūtiet man jūsu izvietojumu, lai varētu plānot mērķa pozīcijas.} \\
            }
        }
    }
     \\ \\

    \theutterance \stepcounter{utterance}  
    & & & \multicolumn{2}{p{0.3\linewidth}}{
        \cellcolor[rgb]{0.9,0.9,0.9}{
            \makecell[{{p{\linewidth}}}]{
                \texttt{\tiny{[GM$|$GM]}}
                \texttt{SAKI: Man pašlaik ir šāds objektu izvietojums: C ir (1,2), L ir (7,5), P ir (7,6). Lūdzu, nosūtiet man jūsu izvietojumu, lai varētu plānot mērķa pozīcijas.} \\
            }
        }
    }
    & & \\ \\

    \theutterance \stepcounter{utterance}  
    & & \multicolumn{4}{p{0.6\linewidth}}{
        \cellcolor[rgb]{0.9,0.9,0.9}{
            \makecell[{{p{\linewidth}}}]{
                \texttt{\tiny{[P1$\langle$GM]}}
                \texttt{Jūsu ziņojums ir nodots otram spēlētājam.} \\
\\ 
\\ 
\texttt{Pašlaik jūs spēlējat kārtu 1 no maksimāli 12.} \\
\texttt{Jūs kopā esat uzkrājuši 0 no 8 sodiem.} \\
\texttt{Otrs spēlētājs nosūtīja šo ziņojumu:} \\
\\ 
\\ 
\texttt{"Man pašlaik ir šāds objektu izvietojums: C ir (1,2), L ir (7,5), P ir (7,6). Lūdzu, nosūtiet man jūsu izvietojumu, lai varētu plānot mērķa pozīcijas."} \\
\\ 
\\ 
\texttt{Kāda ir jūsu nākamā komanda?} \\
            }
        }
    }
    & \\ \\

    \theutterance \stepcounter{utterance}  
    & \multicolumn{4}{p{0.6\linewidth}}{
        \cellcolor[rgb]{0.9,0.9,0.9}{
            \makecell[{{p{\linewidth}}}]{
                \texttt{\tiny{[P1$\rangle$GM]}}
                \texttt{SAKI: Man pašlaik ir šāds objektu izvietojums: C ir (6,3), L ir (5,4), P ir (4,7). Uzskatu, ka mērķa konfigurācija varētu būt šāda: C (1,3), L (2,5), P (3,6). Lūdzu, apstipriniet vai piekļūstiet šim izvietojumam un sagatavojiet gājienus koordinātēm.} \\
            }
        }
    }
    & & \\ \\

    \theutterance \stepcounter{utterance}  
    & & & \multicolumn{2}{p{0.3\linewidth}}{
        \cellcolor[rgb]{0.9,0.9,0.9}{
            \makecell[{{p{\linewidth}}}]{
                \texttt{\tiny{[GM$|$GM]}}
                \texttt{SAKI: Man pašlaik ir šāds objektu izvietojums: C ir (6,3), L ir (5,4), P ir (4,7). Uzskatu, ka mērķa konfigurācija varētu būt šāda: C (1,3), L (2,5), P (3,6). Lūdzu, apstipriniet vai piekļūstiet šim izvietojumam un sagatavojiet gājienus koordinātēm.} \\
            }
        }
    }
    & & \\ \\

    \theutterance \stepcounter{utterance}  
    & & \multicolumn{4}{p{0.6\linewidth}}{
        \cellcolor[rgb]{0.9,0.9,0.9}{
            \makecell[{{p{\linewidth}}}]{
                \texttt{\tiny{[GM$\rangle$P2]}}
                \texttt{Jūsu ziņojums ir nodots otram spēlētājam.} \\
\\ 
\\ 
\texttt{Pašlaik jūs spēlējat kārtu 2 no maksimāli 12.} \\
\texttt{Jūs kopā esat uzkrājuši 0 no 8 sodiem.} \\
\texttt{Otrs spēlētājs nosūtīja šo ziņojumu:} \\
\\ 
\\ 
\texttt{"Man pašlaik ir šāds objektu izvietojums: C ir (6,3), L ir (5,4), P ir (4,7). Uzskatu, ka mērķa konfigurācija varētu būt šāda: C (1,3), L (2,5), P (3,6). Lūdzu, apstipriniet vai piekļūstiet šim izvietojumam un sagatavojiet gājienus koordinātēm."} \\
\\ 
\\ 
\texttt{Kāda ir jūsu nākamā komanda?} \\
            }
        }
    }
    & \\ \\

    \theutterance \stepcounter{utterance}  
    & & & \multicolumn{4}{p{0.6\linewidth}}{
        \cellcolor[rgb]{0.9,0.9,0.9}{
            \makecell[{{p{\linewidth}}}]{
                \texttt{\tiny{[GM$\langle$P2]}}
                \texttt{SAKI: Piekrītu mērķa konfigurācijai: C (1,3), L (2,5), P (3,6). Es gatavoju gājienus. Pārvietošu C uz (1,3), L uz (2,5) un P uz (3,6). Lai apstiprinātu, ir nepieciešams viens no šiem gājieniem: GĀJIENS: C, (1,3); GĀJIENS: L, (2,5); GĀJIENS: P, (3,6) vai visi vienlaikus?} \\
            }
        }
    }
     \\ \\

    \theutterance \stepcounter{utterance}  
    & & & \multicolumn{2}{p{0.3\linewidth}}{
        \cellcolor[rgb]{0.9,0.9,0.9}{
            \makecell[{{p{\linewidth}}}]{
                \texttt{\tiny{[GM$|$GM]}}
                \texttt{SAKI: Piekrītu mērķa konfigurācijai: C (1,3), L (2,5), P (3,6). Es gatavoju gājienus. Pārvietošu C uz (1,3), L uz (2,5) un P uz (3,6). Lai apstiprinātu, ir nepieciešams viens no šiem gājieniem: GĀJIENS: C, (1,3); GĀJIENS: L, (2,5); GĀJIENS: P, (3,6) vai visi vienlaikus?} \\
            }
        }
    }
    & & \\ \\

    \theutterance \stepcounter{utterance}  
    & & & \multicolumn{2}{p{0.3\linewidth}}{
        \cellcolor[rgb]{0.9,0.9,0.9}{
            \makecell[{{p{\linewidth}}}]{
                \texttt{\tiny{[GM$|$GM]}}
                \texttt{Invalid response format: SAKI: Piekrītu mērķa konfigurācijai: C (1,3), L (2,5), P (3,6). Es gatavoju gājienus. Pārvietošu C uz (1,3), L uz (2,5) un P uz (3,6). Lai apstiprinātu, ir nepieciešams viens no šiem gājieniem: GĀJIENS: C, (1,3); GĀJIENS: L, (2,5); GĀJIENS: P, (3,6) vai visi vienlaikus?} \\
            }
        }
    }
    & & \\ \\

    \theutterance \stepcounter{utterance}  
    & & \multicolumn{4}{p{0.6\linewidth}}{
        \cellcolor[rgb]{0.9,0.9,0.9}{
            \makecell[{{p{\linewidth}}}]{
                \texttt{\tiny{[GM$\rangle$P2]}}
                \texttt{Sods: Jūsu ziņojumā ir vairāk nekā viena komanda!} \\
\\ 
\texttt{Pārliecinieties, ka jūsu atbildē ir tikai `SAKI: <ZIŅA>` vai `GĀJIENS: <OBJEKTS>, (<X>, <Y>)`, un nekas cits!} \\
\\ 
\texttt{Jūs kopā esat uzkrājuši 1 no 8 sodiem.} \\
\texttt{Lūdzu, mēģiniet vēlreiz!} \\
            }
        }
    }
    & \\ \\

    \theutterance \stepcounter{utterance}  
    & & & \multicolumn{4}{p{0.6\linewidth}}{
        \cellcolor[rgb]{0.9,0.9,0.9}{
            \makecell[{{p{\linewidth}}}]{
                \texttt{\tiny{[GM$\langle$P2]}}
                \texttt{GĀJIENS: C, (1,3)} \\
            }
        }
    }
     \\ \\

    \theutterance \stepcounter{utterance}  
    & & & \multicolumn{2}{p{0.3\linewidth}}{
        \cellcolor[rgb]{0.9,0.9,0.9}{
            \makecell[{{p{\linewidth}}}]{
                \texttt{\tiny{[GM$|$GM]}}
                \texttt{GĀJIENS: C, (1,3)} \\
            }
        }
    }
    & & \\ \\

    \theutterance \stepcounter{utterance}  
    & & & \multicolumn{2}{p{0.3\linewidth}}{
        \cellcolor[rgb]{0.9,0.9,0.9}{
            \makecell[{{p{\linewidth}}}]{
                \texttt{\tiny{[GM$|$GM]}}
                \texttt{C moved to (1, 3)} \\
            }
        }
    }
    & & \\ \\

    \theutterance \stepcounter{utterance}  
    & & \multicolumn{4}{p{0.6\linewidth}}{
        \cellcolor[rgb]{0.9,0.9,0.9}{
            \makecell[{{p{\linewidth}}}]{
                \texttt{\tiny{[P1$\langle$GM]}}
                \texttt{Jūsu ziņojums ir nodots otram spēlētājam.} \\
\\ 
\\ 
\texttt{Pašlaik jūs spēlējat kārtu 2 no maksimāli 12.} \\
\texttt{Jūs kopā esat uzkrājuši 1 no 8 sodiem.} \\
\texttt{Otrs spēlētājs pārvietoja objektu savā režģī.} \\
\\ 
\\ 
\texttt{Kāda ir jūsu nākamā komanda?} \\
            }
        }
    }
    & \\ \\

    \theutterance \stepcounter{utterance}  
    & \multicolumn{4}{p{0.6\linewidth}}{
        \cellcolor[rgb]{0.9,0.9,0.9}{
            \makecell[{{p{\linewidth}}}]{
                \texttt{\tiny{[P1$\rangle$GM]}}
                \texttt{GĀJIENS: C, (1,3)} \\
            }
        }
    }
    & & \\ \\

    \theutterance \stepcounter{utterance}  
    & & & \multicolumn{2}{p{0.3\linewidth}}{
        \cellcolor[rgb]{0.9,0.9,0.9}{
            \makecell[{{p{\linewidth}}}]{
                \texttt{\tiny{[GM$|$GM]}}
                \texttt{GĀJIENS: C, (1,3)} \\
            }
        }
    }
    & & \\ \\

    \theutterance \stepcounter{utterance}  
    & & & \multicolumn{2}{p{0.3\linewidth}}{
        \cellcolor[rgb]{0.9,0.9,0.9}{
            \makecell[{{p{\linewidth}}}]{
                \texttt{\tiny{[GM$|$GM]}}
                \texttt{C moved to (1, 3)} \\
            }
        }
    }
    & & \\ \\

    \theutterance \stepcounter{utterance}  
    & & \multicolumn{4}{p{0.6\linewidth}}{
        \cellcolor[rgb]{0.9,0.9,0.9}{
            \makecell[{{p{\linewidth}}}]{
                \texttt{\tiny{[GM$\rangle$P2]}}
                \texttt{Veiksmīgi pārvietots 'C' uz (1,3). Jūsu atjauninātais režģis izskatās šādi:} \\
\texttt{```} \\
\texttt{ 1234567} \\
\texttt{╔═╤═══╤═╗} \\
\texttt{║◌│◌◌◌│◌║ 1} \\
\texttt{╟─┼──┬┘◌║ 2} \\
\texttt{║C│◌◌│◌◌║ 3} \\
\texttt{╟─┴──┘◌◌║ 4} \\
\texttt{║◌◌◌◌◌◌◌║ 5} \\
\texttt{║◌◌◌◌◌◌◌║ 6} \\
\texttt{║◌◌◌P◌◌L║ 7} \\
\texttt{╚═══════╝} \\
\texttt{```} \\
\texttt{Pašlaik jūs spēlējat kārtu 3 no maksimāli 12.} \\
\texttt{Jūs kopā esat uzkrājuši 1 no 8 sodiem.} \\
\texttt{Otrs spēlētājs pārvietoja objektu savā režģī.} \\
\\ 
\\ 
\texttt{Kāda ir jūsu nākamā komanda?} \\
            }
        }
    }
    & \\ \\

    \theutterance \stepcounter{utterance}  
    & & & \multicolumn{4}{p{0.6\linewidth}}{
        \cellcolor[rgb]{0.9,0.9,0.9}{
            \makecell[{{p{\linewidth}}}]{
                \texttt{\tiny{[GM$\langle$P2]}}
                \texttt{GĀJIENS: L, (2,5)} \\
            }
        }
    }
     \\ \\

    \theutterance \stepcounter{utterance}  
    & & & \multicolumn{2}{p{0.3\linewidth}}{
        \cellcolor[rgb]{0.9,0.9,0.9}{
            \makecell[{{p{\linewidth}}}]{
                \texttt{\tiny{[GM$|$GM]}}
                \texttt{GĀJIENS: L, (2,5)} \\
            }
        }
    }
    & & \\ \\

    \theutterance \stepcounter{utterance}  
    & & & \multicolumn{2}{p{0.3\linewidth}}{
        \cellcolor[rgb]{0.9,0.9,0.9}{
            \makecell[{{p{\linewidth}}}]{
                \texttt{\tiny{[GM$|$GM]}}
                \texttt{L moved to (2, 5)} \\
            }
        }
    }
    & & \\ \\

    \theutterance \stepcounter{utterance}  
    & & \multicolumn{4}{p{0.6\linewidth}}{
        \cellcolor[rgb]{0.9,0.9,0.9}{
            \makecell[{{p{\linewidth}}}]{
                \texttt{\tiny{[P1$\langle$GM]}}
                \texttt{Veiksmīgi pārvietots 'C' uz (1,3). Jūsu atjauninātais režģis izskatās šādi:} \\
\texttt{```} \\
\texttt{ 1234567} \\
\texttt{╔═╤═══╤═╗} \\
\texttt{║◌│◌◌◌│◌║ 1} \\
\texttt{╟─┼──┬┘◌║ 2} \\
\texttt{║C│◌◌│◌P║ 3} \\
\texttt{╟─┴──┘◌◌║ 4} \\
\texttt{║◌L◌◌◌◌◌║ 5} \\
\texttt{║◌◌◌◌◌◌◌║ 6} \\
\texttt{║◌◌◌◌◌◌◌║ 7} \\
\texttt{╚═══════╝} \\
\texttt{```} \\
\texttt{Pašlaik jūs spēlējat kārtu 3 no maksimāli 12.} \\
\texttt{Jūs kopā esat uzkrājuši 1 no 8 sodiem.} \\
\texttt{Otrs spēlētājs pārvietoja objektu savā režģī.} \\
\\ 
\\ 
\texttt{Kāda ir jūsu nākamā komanda?} \\
            }
        }
    }
    & \\ \\

    \theutterance \stepcounter{utterance}  
    & \multicolumn{4}{p{0.6\linewidth}}{
        \cellcolor[rgb]{0.9,0.9,0.9}{
            \makecell[{{p{\linewidth}}}]{
                \texttt{\tiny{[P1$\rangle$GM]}}
                \texttt{SAKI: Lūdzu, nosūtiet man jūsu pašreizējo objektu izvietojumu, lai varētu plānot nākamos gājiens koordinātes.} \\
            }
        }
    }
    & & \\ \\

    \theutterance \stepcounter{utterance}  
    & & & \multicolumn{2}{p{0.3\linewidth}}{
        \cellcolor[rgb]{0.9,0.9,0.9}{
            \makecell[{{p{\linewidth}}}]{
                \texttt{\tiny{[GM$|$GM]}}
                \texttt{SAKI: Lūdzu, nosūtiet man jūsu pašreizējo objektu izvietojumu, lai varētu plānot nākamos gājiens koordinātes.} \\
            }
        }
    }
    & & \\ \\

    \theutterance \stepcounter{utterance}  
    & & \multicolumn{4}{p{0.6\linewidth}}{
        \cellcolor[rgb]{0.9,0.9,0.9}{
            \makecell[{{p{\linewidth}}}]{
                \texttt{\tiny{[GM$\rangle$P2]}}
                \texttt{Veiksmīgi pārvietots 'L' uz (2,5). Jūsu atjauninātais režģis izskatās šādi:} \\
\texttt{```} \\
\texttt{ 1234567} \\
\texttt{╔═╤═══╤═╗} \\
\texttt{║◌│◌◌◌│◌║ 1} \\
\texttt{╟─┼──┬┘◌║ 2} \\
\texttt{║C│◌◌│◌◌║ 3} \\
\texttt{╟─┴──┘◌◌║ 4} \\
\texttt{║◌L◌◌◌◌◌║ 5} \\
\texttt{║◌◌◌◌◌◌◌║ 6} \\
\texttt{║◌◌◌P◌◌◌║ 7} \\
\texttt{╚═══════╝} \\
\texttt{```} \\
\texttt{Pašlaik jūs spēlējat kārtu 4 no maksimāli 12.} \\
\texttt{Jūs kopā esat uzkrājuši 1 no 8 sodiem.} \\
\texttt{Otrs spēlētājs nosūtīja šo ziņojumu:} \\
\\ 
\\ 
\texttt{"Lūdzu, nosūtiet man jūsu pašreizējo objektu izvietojumu, lai varētu plānot nākamos gājiens koordinātes."} \\
\\ 
\\ 
\texttt{Kāda ir jūsu nākamā komanda?} \\
            }
        }
    }
    & \\ \\

    \theutterance \stepcounter{utterance}  
    & & & \multicolumn{4}{p{0.6\linewidth}}{
        \cellcolor[rgb]{0.9,0.9,0.9}{
            \makecell[{{p{\linewidth}}}]{
                \texttt{\tiny{[GM$\langle$P2]}}
                \texttt{SAKI: Pašlaik ir šāds objektu izvietojums: C (1,2), L (2,5), P (7,6).} \\
            }
        }
    }
     \\ \\

    \theutterance \stepcounter{utterance}  
    & & & \multicolumn{2}{p{0.3\linewidth}}{
        \cellcolor[rgb]{0.9,0.9,0.9}{
            \makecell[{{p{\linewidth}}}]{
                \texttt{\tiny{[GM$|$GM]}}
                \texttt{SAKI: Pašlaik ir šāds objektu izvietojums: C (1,2), L (2,5), P (7,6).} \\
            }
        }
    }
    & & \\ \\

    \theutterance \stepcounter{utterance}  
    & & \multicolumn{4}{p{0.6\linewidth}}{
        \cellcolor[rgb]{0.9,0.9,0.9}{
            \makecell[{{p{\linewidth}}}]{
                \texttt{\tiny{[P1$\langle$GM]}}
                \texttt{Jūsu ziņojums ir nodots otram spēlētājam.} \\
\\ 
\\ 
\texttt{Pašlaik jūs spēlējat kārtu 4 no maksimāli 12.} \\
\texttt{Jūs kopā esat uzkrājuši 1 no 8 sodiem.} \\
\texttt{Otrs spēlētājs nosūtīja šo ziņojumu:} \\
\\ 
\\ 
\texttt{"Pašlaik ir šāds objektu izvietojums: C (1,2), L (2,5), P (7,6)."} \\
\\ 
\\ 
\texttt{Kāda ir jūsu nākamā komanda?} \\
            }
        }
    }
    & \\ \\

    \theutterance \stepcounter{utterance}  
    & \multicolumn{4}{p{0.6\linewidth}}{
        \cellcolor[rgb]{0.9,0.9,0.9}{
            \makecell[{{p{\linewidth}}}]{
                \texttt{\tiny{[P1$\rangle$GM]}}
                \texttt{GĀJIENS: L, (5,2)} \\
            }
        }
    }
    & & \\ \\

    \theutterance \stepcounter{utterance}  
    & & & \multicolumn{2}{p{0.3\linewidth}}{
        \cellcolor[rgb]{0.9,0.9,0.9}{
            \makecell[{{p{\linewidth}}}]{
                \texttt{\tiny{[GM$|$GM]}}
                \texttt{GĀJIENS: L, (5,2)} \\
            }
        }
    }
    & & \\ \\

    \theutterance \stepcounter{utterance}  
    & & & \multicolumn{2}{p{0.3\linewidth}}{
        \cellcolor[rgb]{0.9,0.9,0.9}{
            \makecell[{{p{\linewidth}}}]{
                \texttt{\tiny{[GM$|$GM]}}
                \texttt{Sods: (5,2) nav tukšs, bet satur '┬'.} \\
\\ 
\texttt{Jūs kopā esat uzkrājuši 2 no 8 sodiem.} \\
\texttt{Lūdzu, mēģiniet vēlreiz!} \\
            }
        }
    }
    & & \\ \\

    \theutterance \stepcounter{utterance}  
    & & \multicolumn{4}{p{0.6\linewidth}}{
        \cellcolor[rgb]{0.9,0.9,0.9}{
            \makecell[{{p{\linewidth}}}]{
                \texttt{\tiny{[P1$\langle$GM]}}
                \texttt{Sods: (5,2) nav tukšs, bet satur '┬'.} \\
\\ 
\texttt{Jūs kopā esat uzkrājuši 2 no 8 sodiem.} \\
\texttt{Lūdzu, mēģiniet vēlreiz!} \\
            }
        }
    }
    & \\ \\

    \theutterance \stepcounter{utterance}  
    & \multicolumn{4}{p{0.6\linewidth}}{
        \cellcolor[rgb]{0.9,0.9,0.9}{
            \makecell[{{p{\linewidth}}}]{
                \texttt{\tiny{[P1$\rangle$GM]}}
                \texttt{GĀJIENS: L, (3,5)} \\
            }
        }
    }
    & & \\ \\

    \theutterance \stepcounter{utterance}  
    & & & \multicolumn{2}{p{0.3\linewidth}}{
        \cellcolor[rgb]{0.9,0.9,0.9}{
            \makecell[{{p{\linewidth}}}]{
                \texttt{\tiny{[GM$|$GM]}}
                \texttt{GĀJIENS: L, (3,5)} \\
            }
        }
    }
    & & \\ \\

    \theutterance \stepcounter{utterance}  
    & & & \multicolumn{2}{p{0.3\linewidth}}{
        \cellcolor[rgb]{0.9,0.9,0.9}{
            \makecell[{{p{\linewidth}}}]{
                \texttt{\tiny{[GM$|$GM]}}
                \texttt{L moved to (3, 5)} \\
            }
        }
    }
    & & \\ \\

    \theutterance \stepcounter{utterance}  
    & & \multicolumn{4}{p{0.6\linewidth}}{
        \cellcolor[rgb]{0.9,0.9,0.9}{
            \makecell[{{p{\linewidth}}}]{
                \texttt{\tiny{[GM$\rangle$P2]}}
                \texttt{Jūsu ziņojums ir nodots otram spēlētājam.} \\
\\ 
\\ 
\texttt{Pašlaik jūs spēlējat kārtu 5 no maksimāli 12.} \\
\texttt{Jūs kopā esat uzkrājuši 2 no 8 sodiem.} \\
\texttt{Otrs spēlētājs pārvietoja objektu savā režģī.} \\
\\ 
\\ 
\texttt{Kāda ir jūsu nākamā komanda?} \\
            }
        }
    }
    & \\ \\

    \theutterance \stepcounter{utterance}  
    & & & \multicolumn{4}{p{0.6\linewidth}}{
        \cellcolor[rgb]{0.9,0.9,0.9}{
            \makecell[{{p{\linewidth}}}]{
                \texttt{\tiny{[GM$\langle$P2]}}
                \texttt{GĀJIENS: P, (3,6)} \\
            }
        }
    }
     \\ \\

    \theutterance \stepcounter{utterance}  
    & & & \multicolumn{2}{p{0.3\linewidth}}{
        \cellcolor[rgb]{0.9,0.9,0.9}{
            \makecell[{{p{\linewidth}}}]{
                \texttt{\tiny{[GM$|$GM]}}
                \texttt{GĀJIENS: P, (3,6)} \\
            }
        }
    }
    & & \\ \\

    \theutterance \stepcounter{utterance}  
    & & & \multicolumn{2}{p{0.3\linewidth}}{
        \cellcolor[rgb]{0.9,0.9,0.9}{
            \makecell[{{p{\linewidth}}}]{
                \texttt{\tiny{[GM$|$GM]}}
                \texttt{P moved to (3, 6)} \\
            }
        }
    }
    & & \\ \\

    \theutterance \stepcounter{utterance}  
    & & \multicolumn{4}{p{0.6\linewidth}}{
        \cellcolor[rgb]{0.9,0.9,0.9}{
            \makecell[{{p{\linewidth}}}]{
                \texttt{\tiny{[P1$\langle$GM]}}
                \texttt{Veiksmīgi pārvietots 'L' uz (3,5). Jūsu atjauninātais režģis izskatās šādi:} \\
\texttt{```} \\
\texttt{ 1234567} \\
\texttt{╔═╤═══╤═╗} \\
\texttt{║◌│◌◌◌│◌║ 1} \\
\texttt{╟─┼──┬┘◌║ 2} \\
\texttt{║C│◌◌│◌P║ 3} \\
\texttt{╟─┴──┘◌◌║ 4} \\
\texttt{║◌◌L◌◌◌◌║ 5} \\
\texttt{║◌◌◌◌◌◌◌║ 6} \\
\texttt{║◌◌◌◌◌◌◌║ 7} \\
\texttt{╚═══════╝} \\
\texttt{```} \\
\texttt{Pašlaik jūs spēlējat kārtu 5 no maksimāli 12.} \\
\texttt{Jūs kopā esat uzkrājuši 2 no 8 sodiem.} \\
\texttt{Otrs spēlētājs pārvietoja objektu savā režģī.} \\
\\ 
\\ 
\texttt{Kāda ir jūsu nākamā komanda?} \\
            }
        }
    }
    & \\ \\

    \theutterance \stepcounter{utterance}  
    & \multicolumn{4}{p{0.6\linewidth}}{
        \cellcolor[rgb]{0.9,0.9,0.9}{
            \makecell[{{p{\linewidth}}}]{
                \texttt{\tiny{[P1$\rangle$GM]}}
                \texttt{SAKI: Lūdzu, sagatavojiet man savu pašreizējo objektu izvietojumu, lai varētu plānot nākamos gājiens koordinātas.} \\
            }
        }
    }
    & & \\ \\

    \theutterance \stepcounter{utterance}  
    & & & \multicolumn{2}{p{0.3\linewidth}}{
        \cellcolor[rgb]{0.9,0.9,0.9}{
            \makecell[{{p{\linewidth}}}]{
                \texttt{\tiny{[GM$|$GM]}}
                \texttt{SAKI: Lūdzu, sagatavojiet man savu pašreizējo objektu izvietojumu, lai varētu plānot nākamos gājiens koordinātas.} \\
            }
        }
    }
    & & \\ \\

    \theutterance \stepcounter{utterance}  
    & & \multicolumn{4}{p{0.6\linewidth}}{
        \cellcolor[rgb]{0.9,0.9,0.9}{
            \makecell[{{p{\linewidth}}}]{
                \texttt{\tiny{[GM$\rangle$P2]}}
                \texttt{Veiksmīgi pārvietots 'P' uz (3,6). Jūsu atjauninātais režģis izskatās šādi:} \\
\texttt{```} \\
\texttt{ 1234567} \\
\texttt{╔═╤═══╤═╗} \\
\texttt{║◌│◌◌◌│◌║ 1} \\
\texttt{╟─┼──┬┘◌║ 2} \\
\texttt{║C│◌◌│◌◌║ 3} \\
\texttt{╟─┴──┘◌◌║ 4} \\
\texttt{║◌L◌◌◌◌◌║ 5} \\
\texttt{║◌◌P◌◌◌◌║ 6} \\
\texttt{║◌◌◌◌◌◌◌║ 7} \\
\texttt{╚═══════╝} \\
\texttt{```} \\
\texttt{Pašlaik jūs spēlējat kārtu 6 no maksimāli 12.} \\
\texttt{Jūs kopā esat uzkrājuši 2 no 8 sodiem.} \\
\texttt{Otrs spēlētājs nosūtīja šo ziņojumu:} \\
\\ 
\\ 
\texttt{"Lūdzu, sagatavojiet man savu pašreizējo objektu izvietojumu, lai varētu plānot nākamos gājiens koordinātas."} \\
\\ 
\\ 
\texttt{Kāda ir jūsu nākamā komanda?} \\
            }
        }
    }
    & \\ \\

    \theutterance \stepcounter{utterance}  
    & & & \multicolumn{4}{p{0.6\linewidth}}{
        \cellcolor[rgb]{0.9,0.9,0.9}{
            \makecell[{{p{\linewidth}}}]{
                \texttt{\tiny{[GM$\langle$P2]}}
                \texttt{SAKI: Pašlaik ir šāds objektu izvietojums: C (1,2), L (2,5), P (3,6).} \\
            }
        }
    }
     \\ \\

    \theutterance \stepcounter{utterance}  
    & & & \multicolumn{2}{p{0.3\linewidth}}{
        \cellcolor[rgb]{0.9,0.9,0.9}{
            \makecell[{{p{\linewidth}}}]{
                \texttt{\tiny{[GM$|$GM]}}
                \texttt{SAKI: Pašlaik ir šāds objektu izvietojums: C (1,2), L (2,5), P (3,6).} \\
            }
        }
    }
    & & \\ \\

    \theutterance \stepcounter{utterance}  
    & & \multicolumn{4}{p{0.6\linewidth}}{
        \cellcolor[rgb]{0.9,0.9,0.9}{
            \makecell[{{p{\linewidth}}}]{
                \texttt{\tiny{[P1$\langle$GM]}}
                \texttt{Jūsu ziņojums ir nodots otram spēlētājam.} \\
\\ 
\\ 
\texttt{Pašlaik jūs spēlējat kārtu 6 no maksimāli 12.} \\
\texttt{Jūs kopā esat uzkrājuši 2 no 8 sodiem.} \\
\texttt{Otrs spēlētājs nosūtīja šo ziņojumu:} \\
\\ 
\\ 
\texttt{"Pašlaik ir šāds objektu izvietojums: C (1,2), L (2,5), P (3,6)."} \\
\\ 
\\ 
\texttt{Kāda ir jūsu nākamā komanda?} \\
            }
        }
    }
    & \\ \\

    \theutterance \stepcounter{utterance}  
    & \multicolumn{4}{p{0.6\linewidth}}{
        \cellcolor[rgb]{0.9,0.9,0.9}{
            \makecell[{{p{\linewidth}}}]{
                \texttt{\tiny{[P1$\rangle$GM]}}
                \texttt{GĀJIENS: P, (6,3)} \\
            }
        }
    }
    & & \\ \\

    \theutterance \stepcounter{utterance}  
    & & & \multicolumn{2}{p{0.3\linewidth}}{
        \cellcolor[rgb]{0.9,0.9,0.9}{
            \makecell[{{p{\linewidth}}}]{
                \texttt{\tiny{[GM$|$GM]}}
                \texttt{GĀJIENS: P, (6,3)} \\
            }
        }
    }
    & & \\ \\

    \theutterance \stepcounter{utterance}  
    & & & \multicolumn{2}{p{0.3\linewidth}}{
        \cellcolor[rgb]{0.9,0.9,0.9}{
            \makecell[{{p{\linewidth}}}]{
                \texttt{\tiny{[GM$|$GM]}}
                \texttt{P moved to (6, 3)} \\
            }
        }
    }
    & & \\ \\

    \theutterance \stepcounter{utterance}  
    & & \multicolumn{4}{p{0.6\linewidth}}{
        \cellcolor[rgb]{0.9,0.9,0.9}{
            \makecell[{{p{\linewidth}}}]{
                \texttt{\tiny{[GM$\rangle$P2]}}
                \texttt{Jūsu ziņojums ir nodots otram spēlētājam.} \\
\\ 
\\ 
\texttt{Pašlaik jūs spēlējat kārtu 7 no maksimāli 12.} \\
\texttt{Jūs kopā esat uzkrājuši 2 no 8 sodiem.} \\
\texttt{Otrs spēlētājs pārvietoja objektu savā režģī.} \\
\\ 
\\ 
\texttt{Kāda ir jūsu nākamā komanda?} \\
            }
        }
    }
    & \\ \\

    \theutterance \stepcounter{utterance}  
    & & & \multicolumn{4}{p{0.6\linewidth}}{
        \cellcolor[rgb]{0.9,0.9,0.9}{
            \makecell[{{p{\linewidth}}}]{
                \texttt{\tiny{[GM$\langle$P2]}}
                \texttt{GĀJIENS: C, (1,2)} \\
            }
        }
    }
     \\ \\

    \theutterance \stepcounter{utterance}  
    & & & \multicolumn{2}{p{0.3\linewidth}}{
        \cellcolor[rgb]{0.9,0.9,0.9}{
            \makecell[{{p{\linewidth}}}]{
                \texttt{\tiny{[GM$|$GM]}}
                \texttt{GĀJIENS: C, (1,2)} \\
            }
        }
    }
    & & \\ \\

    \theutterance \stepcounter{utterance}  
    & & & \multicolumn{2}{p{0.3\linewidth}}{
        \cellcolor[rgb]{0.9,0.9,0.9}{
            \makecell[{{p{\linewidth}}}]{
                \texttt{\tiny{[GM$|$GM]}}
                \texttt{Sods: (1,2) nav tukšs, bet satur '─'.} \\
\\ 
\texttt{Jūs kopā esat uzkrājuši 3 no 8 sodiem.} \\
\texttt{Lūdzu, mēģiniet vēlreiz!} \\
            }
        }
    }
    & & \\ \\

    \theutterance \stepcounter{utterance}  
    & & \multicolumn{4}{p{0.6\linewidth}}{
        \cellcolor[rgb]{0.9,0.9,0.9}{
            \makecell[{{p{\linewidth}}}]{
                \texttt{\tiny{[GM$\rangle$P2]}}
                \texttt{Sods: (1,2) nav tukšs, bet satur '─'.} \\
\\ 
\texttt{Jūs kopā esat uzkrājuši 3 no 8 sodiem.} \\
\texttt{Lūdzu, mēģiniet vēlreiz!} \\
            }
        }
    }
    & \\ \\

    \theutterance \stepcounter{utterance}  
    & & & \multicolumn{4}{p{0.6\linewidth}}{
        \cellcolor[rgb]{0.9,0.9,0.9}{
            \makecell[{{p{\linewidth}}}]{
                \texttt{\tiny{[GM$\langle$P2]}}
                \texttt{GĀJIENS: C, (3,5)} \\
            }
        }
    }
     \\ \\

    \theutterance \stepcounter{utterance}  
    & & & \multicolumn{2}{p{0.3\linewidth}}{
        \cellcolor[rgb]{0.9,0.9,0.9}{
            \makecell[{{p{\linewidth}}}]{
                \texttt{\tiny{[GM$|$GM]}}
                \texttt{GĀJIENS: C, (3,5)} \\
            }
        }
    }
    & & \\ \\

    \theutterance \stepcounter{utterance}  
    & & & \multicolumn{2}{p{0.3\linewidth}}{
        \cellcolor[rgb]{0.9,0.9,0.9}{
            \makecell[{{p{\linewidth}}}]{
                \texttt{\tiny{[GM$|$GM]}}
                \texttt{C moved to (3, 5)} \\
            }
        }
    }
    & & \\ \\

    \theutterance \stepcounter{utterance}  
    & & \multicolumn{4}{p{0.6\linewidth}}{
        \cellcolor[rgb]{0.9,0.9,0.9}{
            \makecell[{{p{\linewidth}}}]{
                \texttt{\tiny{[P1$\langle$GM]}}
                \texttt{Veiksmīgi pārvietots 'P' uz (6,3). Jūsu atjauninātais režģis izskatās šādi:} \\
\texttt{```} \\
\texttt{ 1234567} \\
\texttt{╔═╤═══╤═╗} \\
\texttt{║◌│◌◌◌│◌║ 1} \\
\texttt{╟─┼──┬┘◌║ 2} \\
\texttt{║C│◌◌│P◌║ 3} \\
\texttt{╟─┴──┘◌◌║ 4} \\
\texttt{║◌◌L◌◌◌◌║ 5} \\
\texttt{║◌◌◌◌◌◌◌║ 6} \\
\texttt{║◌◌◌◌◌◌◌║ 7} \\
\texttt{╚═══════╝} \\
\texttt{```} \\
\texttt{Pašlaik jūs spēlējat kārtu 7 no maksimāli 12.} \\
\texttt{Jūs kopā esat uzkrājuši 3 no 8 sodiem.} \\
\texttt{Otrs spēlētājs pārvietoja objektu savā režģī.} \\
\\ 
\\ 
\texttt{Kāda ir jūsu nākamā komanda?} \\
            }
        }
    }
    & \\ \\

    \theutterance \stepcounter{utterance}  
    & \multicolumn{4}{p{0.6\linewidth}}{
        \cellcolor[rgb]{0.9,0.9,0.9}{
            \makecell[{{p{\linewidth}}}]{
                \texttt{\tiny{[P1$\rangle$GM]}}
                \texttt{SAKI: Lūdzu, sagatavojiet man savu pašreizējo objektu izvietojumu, lai varētu plānot nākamos gājiens koordinātas.} \\
            }
        }
    }
    & & \\ \\

    \theutterance \stepcounter{utterance}  
    & & & \multicolumn{2}{p{0.3\linewidth}}{
        \cellcolor[rgb]{0.9,0.9,0.9}{
            \makecell[{{p{\linewidth}}}]{
                \texttt{\tiny{[GM$|$GM]}}
                \texttt{SAKI: Lūdzu, sagatavojiet man savu pašreizējo objektu izvietojumu, lai varētu plānot nākamos gājiens koordinātas.} \\
            }
        }
    }
    & & \\ \\

    \theutterance \stepcounter{utterance}  
    & & \multicolumn{4}{p{0.6\linewidth}}{
        \cellcolor[rgb]{0.9,0.9,0.9}{
            \makecell[{{p{\linewidth}}}]{
                \texttt{\tiny{[GM$\rangle$P2]}}
                \texttt{Veiksmīgi pārvietots 'C' uz (3,5). Jūsu atjauninātais režģis izskatās šādi:} \\
\texttt{```} \\
\texttt{ 1234567} \\
\texttt{╔═╤═══╤═╗} \\
\texttt{║◌│◌◌◌│◌║ 1} \\
\texttt{╟─┼──┬┘◌║ 2} \\
\texttt{║◌│◌◌│◌◌║ 3} \\
\texttt{╟─┴──┘◌◌║ 4} \\
\texttt{║◌LC◌◌◌◌║ 5} \\
\texttt{║◌◌P◌◌◌◌║ 6} \\
\texttt{║◌◌◌◌◌◌◌║ 7} \\
\texttt{╚═══════╝} \\
\texttt{```} \\
\texttt{Pašlaik jūs spēlējat kārtu 8 no maksimāli 12.} \\
\texttt{Jūs kopā esat uzkrājuši 3 no 8 sodiem.} \\
\texttt{Otrs spēlētājs nosūtīja šo ziņojumu:} \\
\\ 
\\ 
\texttt{"Lūdzu, sagatavojiet man savu pašreizējo objektu izvietojumu, lai varētu plānot nākamos gājiens koordinātas."} \\
\\ 
\\ 
\texttt{Kāda ir jūsu nākamā komanda?} \\
            }
        }
    }
    & \\ \\

    \theutterance \stepcounter{utterance}  
    & & & \multicolumn{4}{p{0.6\linewidth}}{
        \cellcolor[rgb]{0.9,0.9,0.9}{
            \makecell[{{p{\linewidth}}}]{
                \texttt{\tiny{[GM$\langle$P2]}}
                \texttt{SAKI: Pašlaik ir šāds objektu izvietojums: C (3,5), L (2,5), P (6,7).} \\
            }
        }
    }
     \\ \\

    \theutterance \stepcounter{utterance}  
    & & & \multicolumn{2}{p{0.3\linewidth}}{
        \cellcolor[rgb]{0.9,0.9,0.9}{
            \makecell[{{p{\linewidth}}}]{
                \texttt{\tiny{[GM$|$GM]}}
                \texttt{SAKI: Pašlaik ir šāds objektu izvietojums: C (3,5), L (2,5), P (6,7).} \\
            }
        }
    }
    & & \\ \\

    \theutterance \stepcounter{utterance}  
    & & \multicolumn{4}{p{0.6\linewidth}}{
        \cellcolor[rgb]{0.9,0.9,0.9}{
            \makecell[{{p{\linewidth}}}]{
                \texttt{\tiny{[P1$\langle$GM]}}
                \texttt{Jūsu ziņojums ir nodots otram spēlētājam.} \\
\\ 
\\ 
\texttt{Pašlaik jūs spēlējat kārtu 8 no maksimāli 12.} \\
\texttt{Jūs kopā esat uzkrājuši 3 no 8 sodiem.} \\
\texttt{Otrs spēlētājs nosūtīja šo ziņojumu:} \\
\\ 
\\ 
\texttt{"Pašlaik ir šāds objektu izvietojums: C (3,5), L (2,5), P (6,7)."} \\
\\ 
\\ 
\texttt{Kāda ir jūsu nākamā komanda?} \\
            }
        }
    }
    & \\ \\

    \theutterance \stepcounter{utterance}  
    & \multicolumn{4}{p{0.6\linewidth}}{
        \cellcolor[rgb]{0.9,0.9,0.9}{
            \makecell[{{p{\linewidth}}}]{
                \texttt{\tiny{[P1$\rangle$GM]}}
                \texttt{GĀJIENS: C, (1,4)} \\
            }
        }
    }
    & & \\ \\

    \theutterance \stepcounter{utterance}  
    & & & \multicolumn{2}{p{0.3\linewidth}}{
        \cellcolor[rgb]{0.9,0.9,0.9}{
            \makecell[{{p{\linewidth}}}]{
                \texttt{\tiny{[GM$|$GM]}}
                \texttt{GĀJIENS: C, (1,4)} \\
            }
        }
    }
    & & \\ \\

    \theutterance \stepcounter{utterance}  
    & & & \multicolumn{2}{p{0.3\linewidth}}{
        \cellcolor[rgb]{0.9,0.9,0.9}{
            \makecell[{{p{\linewidth}}}]{
                \texttt{\tiny{[GM$|$GM]}}
                \texttt{Sods: (1,4) nav tukšs, bet satur '─'.} \\
\\ 
\texttt{Jūs kopā esat uzkrājuši 4 no 8 sodiem.} \\
\texttt{Lūdzu, mēģiniet vēlreiz!} \\
            }
        }
    }
    & & \\ \\

    \theutterance \stepcounter{utterance}  
    & & \multicolumn{4}{p{0.6\linewidth}}{
        \cellcolor[rgb]{0.9,0.9,0.9}{
            \makecell[{{p{\linewidth}}}]{
                \texttt{\tiny{[P1$\langle$GM]}}
                \texttt{Sods: (1,4) nav tukšs, bet satur '─'.} \\
\\ 
\texttt{Jūs kopā esat uzkrājuši 4 no 8 sodiem.} \\
\texttt{Lūdzu, mēģiniet vēlreiz!} \\
            }
        }
    }
    & \\ \\

    \theutterance \stepcounter{utterance}  
    & \multicolumn{4}{p{0.6\linewidth}}{
        \cellcolor[rgb]{0.9,0.9,0.9}{
            \makecell[{{p{\linewidth}}}]{
                \texttt{\tiny{[P1$\rangle$GM]}}
                \texttt{GĀJIENS: C, (2,4)} \\
            }
        }
    }
    & & \\ \\

    \theutterance \stepcounter{utterance}  
    & & & \multicolumn{2}{p{0.3\linewidth}}{
        \cellcolor[rgb]{0.9,0.9,0.9}{
            \makecell[{{p{\linewidth}}}]{
                \texttt{\tiny{[GM$|$GM]}}
                \texttt{GĀJIENS: C, (2,4)} \\
            }
        }
    }
    & & \\ \\

    \theutterance \stepcounter{utterance}  
    & & & \multicolumn{2}{p{0.3\linewidth}}{
        \cellcolor[rgb]{0.9,0.9,0.9}{
            \makecell[{{p{\linewidth}}}]{
                \texttt{\tiny{[GM$|$GM]}}
                \texttt{Sods: (2,4) nav tukšs, bet satur '┴'.} \\
\\ 
\texttt{Jūs kopā esat uzkrājuši 5 no 8 sodiem.} \\
\texttt{Lūdzu, mēģiniet vēlreiz!} \\
            }
        }
    }
    & & \\ \\

    \theutterance \stepcounter{utterance}  
    & & \multicolumn{4}{p{0.6\linewidth}}{
        \cellcolor[rgb]{0.9,0.9,0.9}{
            \makecell[{{p{\linewidth}}}]{
                \texttt{\tiny{[P1$\langle$GM]}}
                \texttt{Sods: (2,4) nav tukšs, bet satur '┴'.} \\
\\ 
\texttt{Jūs kopā esat uzkrājuši 5 no 8 sodiem.} \\
\texttt{Lūdzu, mēģiniet vēlreiz!} \\
            }
        }
    }
    & \\ \\

    \theutterance \stepcounter{utterance}  
    & \multicolumn{4}{p{0.6\linewidth}}{
        \cellcolor[rgb]{0.9,0.9,0.9}{
            \makecell[{{p{\linewidth}}}]{
                \texttt{\tiny{[P1$\rangle$GM]}}
                \texttt{GĀJIENS: C, (4,1)} \\
            }
        }
    }
    & & \\ \\

    \theutterance \stepcounter{utterance}  
    & & & \multicolumn{2}{p{0.3\linewidth}}{
        \cellcolor[rgb]{0.9,0.9,0.9}{
            \makecell[{{p{\linewidth}}}]{
                \texttt{\tiny{[GM$|$GM]}}
                \texttt{GĀJIENS: C, (4,1)} \\
            }
        }
    }
    & & \\ \\

    \theutterance \stepcounter{utterance}  
    & & & \multicolumn{2}{p{0.3\linewidth}}{
        \cellcolor[rgb]{0.9,0.9,0.9}{
            \makecell[{{p{\linewidth}}}]{
                \texttt{\tiny{[GM$|$GM]}}
                \texttt{C moved to (4, 1)} \\
            }
        }
    }
    & & \\ \\

    \theutterance \stepcounter{utterance}  
    & & \multicolumn{4}{p{0.6\linewidth}}{
        \cellcolor[rgb]{0.9,0.9,0.9}{
            \makecell[{{p{\linewidth}}}]{
                \texttt{\tiny{[GM$\rangle$P2]}}
                \texttt{Jūsu ziņojums ir nodots otram spēlētājam.} \\
\\ 
\\ 
\texttt{Pašlaik jūs spēlējat kārtu 9 no maksimāli 12.} \\
\texttt{Jūs kopā esat uzkrājuši 5 no 8 sodiem.} \\
\texttt{Otrs spēlētājs pārvietoja objektu savā režģī.} \\
\\ 
\\ 
\texttt{Kāda ir jūsu nākamā komanda?} \\
            }
        }
    }
    & \\ \\

    \theutterance \stepcounter{utterance}  
    & & & \multicolumn{4}{p{0.6\linewidth}}{
        \cellcolor[rgb]{0.9,0.9,0.9}{
            \makecell[{{p{\linewidth}}}]{
                \texttt{\tiny{[GM$\langle$P2]}}
                \texttt{GĀJIENS: L, (2,6)} \\
            }
        }
    }
     \\ \\

    \theutterance \stepcounter{utterance}  
    & & & \multicolumn{2}{p{0.3\linewidth}}{
        \cellcolor[rgb]{0.9,0.9,0.9}{
            \makecell[{{p{\linewidth}}}]{
                \texttt{\tiny{[GM$|$GM]}}
                \texttt{GĀJIENS: L, (2,6)} \\
            }
        }
    }
    & & \\ \\

    \theutterance \stepcounter{utterance}  
    & & & \multicolumn{2}{p{0.3\linewidth}}{
        \cellcolor[rgb]{0.9,0.9,0.9}{
            \makecell[{{p{\linewidth}}}]{
                \texttt{\tiny{[GM$|$GM]}}
                \texttt{L moved to (2, 6)} \\
            }
        }
    }
    & & \\ \\

    \theutterance \stepcounter{utterance}  
    & & \multicolumn{4}{p{0.6\linewidth}}{
        \cellcolor[rgb]{0.9,0.9,0.9}{
            \makecell[{{p{\linewidth}}}]{
                \texttt{\tiny{[P1$\langle$GM]}}
                \texttt{Veiksmīgi pārvietots 'C' uz (4,1). Jūsu atjauninātais režģis izskatās šādi:} \\
\texttt{```} \\
\texttt{ 1234567} \\
\texttt{╔═╤═══╤═╗} \\
\texttt{║◌│◌C◌│◌║ 1} \\
\texttt{╟─┼──┬┘◌║ 2} \\
\texttt{║◌│◌◌│P◌║ 3} \\
\texttt{╟─┴──┘◌◌║ 4} \\
\texttt{║◌◌L◌◌◌◌║ 5} \\
\texttt{║◌◌◌◌◌◌◌║ 6} \\
\texttt{║◌◌◌◌◌◌◌║ 7} \\
\texttt{╚═══════╝} \\
\texttt{```} \\
\texttt{Pašlaik jūs spēlējat kārtu 9 no maksimāli 12.} \\
\texttt{Jūs kopā esat uzkrājuši 5 no 8 sodiem.} \\
\texttt{Otrs spēlētājs pārvietoja objektu savā režģī.} \\
\\ 
\\ 
\texttt{Kāda ir jūsu nākamā komanda?} \\
            }
        }
    }
    & \\ \\

    \theutterance \stepcounter{utterance}  
    & \multicolumn{4}{p{0.6\linewidth}}{
        \cellcolor[rgb]{0.9,0.9,0.9}{
            \makecell[{{p{\linewidth}}}]{
                \texttt{\tiny{[P1$\rangle$GM]}}
                \texttt{SAKI: Lūdzu, sagatavojiet man savu pašreizējo objektu izvietojumu, lai varētu plānot nākamos gājiens koordinātas.} \\
            }
        }
    }
    & & \\ \\

    \theutterance \stepcounter{utterance}  
    & & & \multicolumn{2}{p{0.3\linewidth}}{
        \cellcolor[rgb]{0.9,0.9,0.9}{
            \makecell[{{p{\linewidth}}}]{
                \texttt{\tiny{[GM$|$GM]}}
                \texttt{SAKI: Lūdzu, sagatavojiet man savu pašreizējo objektu izvietojumu, lai varētu plānot nākamos gājiens koordinātas.} \\
            }
        }
    }
    & & \\ \\

    \theutterance \stepcounter{utterance}  
    & & \multicolumn{4}{p{0.6\linewidth}}{
        \cellcolor[rgb]{0.9,0.9,0.9}{
            \makecell[{{p{\linewidth}}}]{
                \texttt{\tiny{[GM$\rangle$P2]}}
                \texttt{Veiksmīgi pārvietots 'L' uz (2,6). Jūsu atjauninātais režģis izskatās šādi:} \\
\texttt{```} \\
\texttt{ 1234567} \\
\texttt{╔═╤═══╤═╗} \\
\texttt{║◌│◌◌◌│◌║ 1} \\
\texttt{╟─┼──┬┘◌║ 2} \\
\texttt{║◌│◌◌│◌◌║ 3} \\
\texttt{╟─┴──┘◌◌║ 4} \\
\texttt{║◌◌C◌◌◌◌║ 5} \\
\texttt{║◌LP◌◌◌◌║ 6} \\
\texttt{║◌◌◌◌◌◌◌║ 7} \\
\texttt{╚═══════╝} \\
\texttt{```} \\
\texttt{Pašlaik jūs spēlējat kārtu 10 no maksimāli 12.} \\
\texttt{Jūs kopā esat uzkrājuši 5 no 8 sodiem.} \\
\texttt{Otrs spēlētājs nosūtīja šo ziņojumu:} \\
\\ 
\\ 
\texttt{"Lūdzu, sagatavojiet man savu pašreizējo objektu izvietojumu, lai varētu plānot nākamos gājiens koordinātas."} \\
\\ 
\\ 
\texttt{Kāda ir jūsu nākamā komanda?} \\
            }
        }
    }
    & \\ \\

    \theutterance \stepcounter{utterance}  
    & & & \multicolumn{4}{p{0.6\linewidth}}{
        \cellcolor[rgb]{0.9,0.9,0.9}{
            \makecell[{{p{\linewidth}}}]{
                \texttt{\tiny{[GM$\langle$P2]}}
                \texttt{SAKI: Pašlaik ir šāds objektu izvietojums: C (5,2), L (2,6), P (6,7).} \\
            }
        }
    }
     \\ \\

    \theutterance \stepcounter{utterance}  
    & & & \multicolumn{2}{p{0.3\linewidth}}{
        \cellcolor[rgb]{0.9,0.9,0.9}{
            \makecell[{{p{\linewidth}}}]{
                \texttt{\tiny{[GM$|$GM]}}
                \texttt{SAKI: Pašlaik ir šāds objektu izvietojums: C (5,2), L (2,6), P (6,7).} \\
            }
        }
    }
    & & \\ \\

    \theutterance \stepcounter{utterance}  
    & & \multicolumn{4}{p{0.6\linewidth}}{
        \cellcolor[rgb]{0.9,0.9,0.9}{
            \makecell[{{p{\linewidth}}}]{
                \texttt{\tiny{[P1$\langle$GM]}}
                \texttt{Jūsu ziņojums ir nodots otram spēlētājam.} \\
\\ 
\\ 
\texttt{Pašlaik jūs spēlējat kārtu 10 no maksimāli 12.} \\
\texttt{Jūs kopā esat uzkrājuši 5 no 8 sodiem.} \\
\texttt{Otrs spēlētājs nosūtīja šo ziņojumu:} \\
\\ 
\\ 
\texttt{"Pašlaik ir šāds objektu izvietojums: C (5,2), L (2,6), P (6,7)."} \\
\\ 
\\ 
\texttt{Kāda ir jūsu nākamā komanda?} \\
            }
        }
    }
    & \\ \\

    \theutterance \stepcounter{utterance}  
    & \multicolumn{4}{p{0.6\linewidth}}{
        \cellcolor[rgb]{0.9,0.9,0.9}{
            \makecell[{{p{\linewidth}}}]{
                \texttt{\tiny{[P1$\rangle$GM]}}
                \texttt{GĀJIENS: L, (4,6)} \\
            }
        }
    }
    & & \\ \\

    \theutterance \stepcounter{utterance}  
    & & & \multicolumn{2}{p{0.3\linewidth}}{
        \cellcolor[rgb]{0.9,0.9,0.9}{
            \makecell[{{p{\linewidth}}}]{
                \texttt{\tiny{[GM$|$GM]}}
                \texttt{GĀJIENS: L, (4,6)} \\
            }
        }
    }
    & & \\ \\

    \theutterance \stepcounter{utterance}  
    & & & \multicolumn{2}{p{0.3\linewidth}}{
        \cellcolor[rgb]{0.9,0.9,0.9}{
            \makecell[{{p{\linewidth}}}]{
                \texttt{\tiny{[GM$|$GM]}}
                \texttt{L moved to (4, 6)} \\
            }
        }
    }
    & & \\ \\

    \theutterance \stepcounter{utterance}  
    & & \multicolumn{4}{p{0.6\linewidth}}{
        \cellcolor[rgb]{0.9,0.9,0.9}{
            \makecell[{{p{\linewidth}}}]{
                \texttt{\tiny{[GM$\rangle$P2]}}
                \texttt{Jūsu ziņojums ir nodots otram spēlētājam.} \\
\\ 
\\ 
\texttt{Pašlaik jūs spēlējat kārtu 11 no maksimāli 12.} \\
\texttt{Jūs kopā esat uzkrājuši 5 no 8 sodiem.} \\
\texttt{Otrs spēlētājs pārvietoja objektu savā režģī.} \\
\\ 
\\ 
\texttt{Kāda ir jūsu nākamā komanda?} \\
            }
        }
    }
    & \\ \\

    \theutterance \stepcounter{utterance}  
    & & & \multicolumn{4}{p{0.6\linewidth}}{
        \cellcolor[rgb]{0.9,0.9,0.9}{
            \makecell[{{p{\linewidth}}}]{
                \texttt{\tiny{[GM$\langle$P2]}}
                \texttt{GĀJIENS: P, (6,5)} \\
            }
        }
    }
     \\ \\

    \theutterance \stepcounter{utterance}  
    & & & \multicolumn{2}{p{0.3\linewidth}}{
        \cellcolor[rgb]{0.9,0.9,0.9}{
            \makecell[{{p{\linewidth}}}]{
                \texttt{\tiny{[GM$|$GM]}}
                \texttt{GĀJIENS: P, (6,5)} \\
            }
        }
    }
    & & \\ \\

    \theutterance \stepcounter{utterance}  
    & & & \multicolumn{2}{p{0.3\linewidth}}{
        \cellcolor[rgb]{0.9,0.9,0.9}{
            \makecell[{{p{\linewidth}}}]{
                \texttt{\tiny{[GM$|$GM]}}
                \texttt{P moved to (6, 5)} \\
            }
        }
    }
    & & \\ \\

    \theutterance \stepcounter{utterance}  
    & & \multicolumn{4}{p{0.6\linewidth}}{
        \cellcolor[rgb]{0.9,0.9,0.9}{
            \makecell[{{p{\linewidth}}}]{
                \texttt{\tiny{[P1$\langle$GM]}}
                \texttt{Veiksmīgi pārvietots 'L' uz (4,6). Jūsu atjauninātais režģis izskatās šādi:} \\
\texttt{```} \\
\texttt{ 1234567} \\
\texttt{╔═╤═══╤═╗} \\
\texttt{║◌│◌C◌│◌║ 1} \\
\texttt{╟─┼──┬┘◌║ 2} \\
\texttt{║◌│◌◌│P◌║ 3} \\
\texttt{╟─┴──┘◌◌║ 4} \\
\texttt{║◌◌◌◌◌◌◌║ 5} \\
\texttt{║◌◌◌L◌◌◌║ 6} \\
\texttt{║◌◌◌◌◌◌◌║ 7} \\
\texttt{╚═══════╝} \\
\texttt{```} \\
\texttt{Pašlaik jūs spēlējat kārtu 11 no maksimāli 12.} \\
\texttt{Jūs kopā esat uzkrājuši 5 no 8 sodiem.} \\
\texttt{Otrs spēlētājs pārvietoja objektu savā režģī.} \\
\\ 
\\ 
\texttt{Kāda ir jūsu nākamā komanda?} \\
            }
        }
    }
    & \\ \\

    \theutterance \stepcounter{utterance}  
    & \multicolumn{4}{p{0.6\linewidth}}{
        \cellcolor[rgb]{0.9,0.9,0.9}{
            \makecell[{{p{\linewidth}}}]{
                \texttt{\tiny{[P1$\rangle$GM]}}
                \texttt{SAKI: Lūdzu, sagatavojiet man jūsu pašreizējo objektu izvietojumu, lai varētu nolēmt mērķa pozīcijas pirmajiem trim objektiem, kas atbilst mērķa konfigurācijai.} \\
            }
        }
    }
    & & \\ \\

    \theutterance \stepcounter{utterance}  
    & & & \multicolumn{2}{p{0.3\linewidth}}{
        \cellcolor[rgb]{0.9,0.9,0.9}{
            \makecell[{{p{\linewidth}}}]{
                \texttt{\tiny{[GM$|$GM]}}
                \texttt{SAKI: Lūdzu, sagatavojiet man jūsu pašreizējo objektu izvietojumu, lai varētu nolēmt mērķa pozīcijas pirmajiem trim objektiem, kas atbilst mērķa konfigurācijai.} \\
            }
        }
    }
    & & \\ \\

    \theutterance \stepcounter{utterance}  
    & & \multicolumn{4}{p{0.6\linewidth}}{
        \cellcolor[rgb]{0.9,0.9,0.9}{
            \makecell[{{p{\linewidth}}}]{
                \texttt{\tiny{[GM$\rangle$P2]}}
                \texttt{Veiksmīgi pārvietots 'P' uz (6,5). Jūsu atjauninātais režģis izskatās šādi:} \\
\texttt{```} \\
\texttt{ 1234567} \\
\texttt{╔═╤═══╤═╗} \\
\texttt{║◌│◌◌◌│◌║ 1} \\
\texttt{╟─┼──┬┘◌║ 2} \\
\texttt{║◌│◌◌│◌◌║ 3} \\
\texttt{╟─┴──┘◌◌║ 4} \\
\texttt{║◌◌C◌◌P◌║ 5} \\
\texttt{║◌L◌◌◌◌◌║ 6} \\
\texttt{║◌◌◌◌◌◌◌║ 7} \\
\texttt{╚═══════╝} \\
\texttt{```} \\
\texttt{Pašlaik jūs spēlējat kārtu 12 no maksimāli 12.} \\
\texttt{Jūs kopā esat uzkrājuši 5 no 8 sodiem.} \\
\texttt{Otrs spēlētājs nosūtīja šo ziņojumu:} \\
\\ 
\\ 
\texttt{"Lūdzu, sagatavojiet man jūsu pašreizējo objektu izvietojumu, lai varētu nolēmt mērķa pozīcijas pirmajiem trim objektiem, kas atbilst mērķa konfigurācijai."} \\
\\ 
\\ 
\texttt{Kāda ir jūsu nākamā komanda?} \\
            }
        }
    }
    & \\ \\

    \theutterance \stepcounter{utterance}  
    & & & \multicolumn{4}{p{0.6\linewidth}}{
        \cellcolor[rgb]{0.9,0.9,0.9}{
            \makecell[{{p{\linewidth}}}]{
                \texttt{\tiny{[GM$\langle$P2]}}
                \texttt{SAKI: Pašlaik ir šāds objektu izvietojums: C (5,2), L (6,3), P (5,6).} \\
            }
        }
    }
     \\ \\

    \theutterance \stepcounter{utterance}  
    & & & \multicolumn{2}{p{0.3\linewidth}}{
        \cellcolor[rgb]{0.9,0.9,0.9}{
            \makecell[{{p{\linewidth}}}]{
                \texttt{\tiny{[GM$|$GM]}}
                \texttt{SAKI: Pašlaik ir šāds objektu izvietojums: C (5,2), L (6,3), P (5,6).} \\
            }
        }
    }
    & & \\ \\

    \theutterance \stepcounter{utterance}  
    & & & \multicolumn{2}{p{0.3\linewidth}}{
        \cellcolor[rgb]{0.9,0.9,0.9}{
            \makecell[{{p{\linewidth}}}]{
                \texttt{\tiny{[GM$|$GM]}}
                \texttt{Maximum number of rounds reached} \\
            }
        }
    }
    & & \\ \\

    \theutterance \stepcounter{utterance}  
    & & & \multicolumn{2}{p{0.3\linewidth}}{
        \cellcolor[rgb]{0.9,0.9,0.9}{
            \makecell[{{p{\linewidth}}}]{
                \texttt{\tiny{[GM$|$GM]}}
                \texttt{Initial states:} \\
\texttt{```} \\
\texttt{Player 1:} \\
\texttt{ 1234567} \\
\texttt{╔═╤═══╤═╗} \\
\texttt{║◌│◌◌◌│◌║ 1} \\
\texttt{╟─┼──┬┘◌║ 2} \\
\texttt{║◌│◌◌│◌P║ 3} \\
\texttt{╟─┴──┘◌◌║ 4} \\
\texttt{║◌L◌◌◌◌◌║ 5} \\
\texttt{║◌C◌◌◌◌◌║ 6} \\
\texttt{║◌◌◌◌◌◌◌║ 7} \\
\texttt{╚═══════╝} \\
\\ 
\\ 
\texttt{Player 2:} \\
\texttt{ 1234567} \\
\texttt{╔═╤═══╤═╗} \\
\texttt{║◌│◌C◌│◌║ 1} \\
\texttt{╟─┼──┬┘◌║ 2} \\
\texttt{║◌│◌◌│◌◌║ 3} \\
\texttt{╟─┴──┘◌◌║ 4} \\
\texttt{║◌◌◌◌◌◌◌║ 5} \\
\texttt{║◌◌◌◌◌◌◌║ 6} \\
\texttt{║◌◌◌P◌◌L║ 7} \\
\texttt{╚═══════╝} \\
\\ 
\texttt{```} \\
            }
        }
    }
    & & \\ \\

    \theutterance \stepcounter{utterance}  
    & & & \multicolumn{2}{p{0.3\linewidth}}{
        \cellcolor[rgb]{0.9,0.9,0.9}{
            \makecell[{{p{\linewidth}}}]{
                \texttt{\tiny{[GM$|$GM]}}
                \texttt{End states:} \\
\texttt{```} \\
\texttt{Player 1:} \\
\texttt{ 1234567} \\
\texttt{╔═╤═══╤═╗} \\
\texttt{║◌│◌C◌│◌║ 1} \\
\texttt{╟─┼──┬┘◌║ 2} \\
\texttt{║◌│◌◌│P◌║ 3} \\
\texttt{╟─┴──┘◌◌║ 4} \\
\texttt{║◌◌◌◌◌◌◌║ 5} \\
\texttt{║◌◌◌L◌◌◌║ 6} \\
\texttt{║◌◌◌◌◌◌◌║ 7} \\
\texttt{╚═══════╝} \\
\\ 
\\ 
\texttt{Player 2:} \\
\texttt{ 1234567} \\
\texttt{╔═╤═══╤═╗} \\
\texttt{║◌│◌◌◌│◌║ 1} \\
\texttt{╟─┼──┬┘◌║ 2} \\
\texttt{║◌│◌◌│◌◌║ 3} \\
\texttt{╟─┴──┘◌◌║ 4} \\
\texttt{║◌◌C◌◌P◌║ 5} \\
\texttt{║◌L◌◌◌◌◌║ 6} \\
\texttt{║◌◌◌◌◌◌◌║ 7} \\
\texttt{╚═══════╝} \\
\\ 
\texttt{```} \\
            }
        }
    }
    & & \\ \\

    \theutterance \stepcounter{utterance}  
    & & & \multicolumn{2}{p{0.3\linewidth}}{
        \cellcolor[rgb]{0.9,0.9,0.9}{
            \makecell[{{p{\linewidth}}}]{
                \texttt{\tiny{[GM$|$GM]}}
                \texttt{* success: True} \\
\texttt{* lose: False} \\
\texttt{* aborted: False} \\
\texttt{{-}{-}{-}{-}{-}{-}{-}} \\
\texttt{* Shifts: 5.00} \\
\texttt{* Max Shifts: 4.00} \\
\texttt{* Min Shifts: 2.00} \\
\texttt{* End Distance Sum: 8.12} \\
\texttt{* Init Distance Sum: 15.77} \\
\texttt{* Expected Distance Sum: 12.57} \\
\texttt{* Penalties: 5.00} \\
\texttt{* Max Penalties: 8.00} \\
\texttt{* Rounds: 12.00} \\
\texttt{* Max Rounds: 12.00} \\
\texttt{* Object Count: 3.00} \\
            }
        }
    }
    & & \\ \\

\end{supertabular}
}

\end{document}
